\chapter{[KHD] Katastrophen, Gro{\ss}schadensereignisse,
Gefahrengutunfall}
\label{chp:KHD}
\AasRsN{RS~XI}

\emph{Maintainer: Roman Koch}

% Changelog:
% 
% 2009-02-20: Fertig für Kurs 2009/02.
% 
% 2009-02-25: Start Changelog.
% 
%\write18{git log > version.log}

%\input{|"git log"}

% \A{\tiny
% \input{include-khd-chapter.commit20.log} 
% }

 
\minitoc 

\section{Katastrophen, Großschadensereignisse, Unfall}
% TODO : Layout-Anpassung
\begin{anfxwarning}[author=GABS]{Layout-Anpassung notwendig.}
   Dieser Teil entspricht nicht dem Standard-Layout!
\end{anfxwarning}


\begin{minipage}{\linewidth}
\begin{center}

\includegraphics[width=\linewidth]{./images/noauth/AASdev20090228-img109.jpg}
 
\end{center}

\Parallel{
\includegraphics[width=6.212cm,height=3.933cm]{./images/noauth/AASdev20090228-img110.jpg}
}{
\Captionof{figure}{ICE-Unglück von Eschede am 3. Juni
1998. \SourcePicture{Wikipedia, Nils Fretwurst; Lizenz: gemein frei
({\quotedblbase}public domain{\textquotedblleft})}}

Verstorbene: 101

Schwerverletzte: 88  

Leicht- und Unverletzte:106

}

\end{minipage}

\subsubsection{Definition \T{Unfall} (Bagatelle, Notfall)}
% TODO : Layout-Anpassung
\begin{anfxwarning}[author=GABS]{Layout-Anpassung notwendig.}
   Dieser Teil entspricht nicht dem Standard-Layout!
\end{anfxwarning}


\begin{tabulary}{\linewidth}{LL}
\toprule
\TabSub{Ereignis:}
 &
plötzlich, unerwartet
\\\addlinespace
\TabSub{Dauer:}
 &
Minuten bis 1 Std.
\\\addlinespace
\TabSub{Infrastruktur:}\MarginBugfix{\emph{Infrastruktur:} Wichtige
Basiseinrichtungen und Dienstleistungen, die das öffentliche Leben ermöglichen;
z.\,B. Straßen, Elektrizität, Telefon, Krankenhäuser, Verwaltungseinrichtungen, 
Nahrungs- und Wasserversorgung, Abfall- und Abwasserentsorgung, usw.} 
&
intakt 
\\\addlinespace
\TabSub{Verletzte:}
 &
bis 14
\\\addlinespace
\TabSub{Erstversorgung:}
 &
innerhalb von Minuten durch Rettungsdienst nach Regeln der
Individual\-medizin
\\\addlinespace
\TabSub{Abtransport:}
 &
unbehindert Routinebetrieb mit örtlich vorhandenen Mitteln
\\\addlinespace
\TabSub{Definitive Versorgung:}
 &
in umliegende Krankenhäuser nach Regeln der Individualmedizin
\\\addlinespace
\TabSub{Material und Mannschaft:}
 &
Normalbetrieb Rettungs- und Krankentransport
\\\bottomrule
\end{tabulary}


\subsubsection{Definition \T{Gro{\ss}unfall} (Gro{\ss}schadensereignis)}
\index{Gro{\ss}unfall}
\index{Gro{\ss}schadensereignis}
% TODO : Layout-Anpassung
\begin{anfxwarning}[author=GABS]{Layout-Anpassung notwendig.}
   Dieser Teil entspricht nicht dem Standard-Layout!
\end{anfxwarning}
\Definition{{\quotedblbase}Ein Gro{\ss}unfall liegt vor, wenn anzunehmen ist,
dass das Ereignis mit den normalen örtlichen personellen und materiellen
Kräften und Mitteln nicht bewältigt werden kann, aber keine
erklärte Katastrophensituation vorliegt{\textquotedblleft}}


\begin{tabulary}{\linewidth}{LL}
\toprule
\TabSub{Ereignis:}
 &
lokal begrenzt
\\\addlinespace
\TabSub{Dauer:}
 &
mehrere Stunden
\\\addlinespace
\TabSub{Infrastruktur:}
 &
intakt
\\\addlinespace
\TabSub{Verletzte:}
 &
ab 15
\\\addlinespace
\TabSub{Erstversorgung:}
 &
innerhalb von Minuten bis Stunden durch den Rettungsdienst nach den
Regeln der Katastrophenmedizin
\\\addlinespace
\TabSub{Abtransport:}
 &
relativ unbehindert nach der Erstversorgung, genügend vorhandene
Transportmittel
\\\addlinespace
\TabSub{Definitive Versorgung:}
 &
in umliegende Krankenhäuser nach den Regeln der Individualmedizin
\\\addlinespace
\TabSub{Material und Mannschaft:}
 &
alle verfügbaren Helfer, Normalbetrieb Rettungs- und Krankentransport
inkl. Sondereinsatzgruppe und zusätzlichem Material
\\\bottomrule
\end{tabulary}


\subsubsection{Definition \index{Katastrophe}\T{Katastrophe}
(au{\ss}ergewöhnliche Lage)}
% TODO : Layout-Anpassung
\begin{anfxwarning}[author=GABS]{Layout-Anpassung notwendig.}
   Dieser Teil entspricht nicht dem Standard-Layout!
\end{anfxwarning}

\emph{griech.:  {\quotedblbase}Umkehr der normalen
Verhältnisse{\textquotedblleft}} 

\Definition{Ein \Es{unvorhergesehenes Ereignis} mit einem
\Es{au{\ss}ergewöhnlichen Schadensausma{\ss}} das unmittelbar bevorsteht oder
bereits eingetreten ist, eine \Es{konkrete Gefahr für Menschen}\Z{, Tiere,
Umwelt, Kulturgüter} und Sachwerte sowie für die \Es{Infrastruktur} zur
Sicherstellung der Versorgung mit lebensnotwendigen Gütern und Dienstleistungen, welches
der \Es{koordinierten Führung durch die Behörde} bedarf. }

\begin{center}
\begin{tabular}{m{3.31cm}m{9.99cm}}
\toprule
\TabSub{Ereignis:}
 &
(über)-regionales Gebiet
\\\addlinespace
\TabSub{Dauer:}
 &
Tage, Wochen
\\\addlinespace
\TabSub{Infrastruktur:}
 &
gestört, zerstört
\\\addlinespace
\TabSub{Verletzte:}
 &
keine bis viele, je nach Art des Ereignisses
\\\addlinespace
\TabSub{Erstversorgung:}
 &
innerhalb von Stunden bis Tagen durch KHD\footnotemark{}{}-Einheiten
nach den Regeln der Katastrophenmedizin
\\\addlinespace
\TabSub{Abtransport:}
 &
stark behindert verzögert, wenige/nicht vorhandene Transportmittel
\\\addlinespace
\TabSub{Definitive Versorgung:}
 &
auf Verbandsplätzen in Behelfsspitälern, in Lazaretten nach Regeln
der Katastrophenmedizin
\\\addlinespace
\TabSub{Material und Mannschaft:}
 &
\"Uberlebende und verfügbare Helfer, einsatzbereite Fahrzeug,
Kat-Einsatz aller Hilfsorganisationen, überregionale und
internationale Hilfe
\\\bottomrule
\end{tabular}
\end{center}
\footnotetext{KHD: Katastrophenhilfsdienst}

\subsubsection{Definition \T{Individualmedizin}}\index{Individualmedizin}
% TODO : Layout-Anpassung
\begin{anfxwarning}[author=GABS]{Layout-Anpassung notwendig.}
   Dieser Teil entspricht nicht dem Standard-Layout!
\end{anfxwarning}

\Definition{Bedarfsgerechte und bestmögliche Versorgung jedes Patienten.} 
% TODO : Layout-Anpassung
\begin{anfxwarning}[author=GABS]{Layout-Anpassung notwendig.}
   Dieser Teil entspricht nicht dem Standard-Layout!
\end{anfxwarning}

Es besteht freie Arztwahl, welche auf das Vertrauensverhältnis zwischen
Patient und Arzt ausgerichtet ist.

\subsubsection{Definition \T{Katastrophenmedizin}}\index{Katastrophenmedizin}

\Definition{Massenversorgung von Verletzten und
Erkrankten mit beschränkten Mitteln und dem Zwang zur Einteilung nach
Dringlichkeit.}

Eine Vielzahl von Hilfsbedürftigen stehen wenigen \"Arzten,
Sanitätern und sonstigen Helfern mit oft unzureichenden Mitteln
gegenüber. Ziel muss es sein, unter den gegebenen Umständen so
viele Leben wie möglich zu retten, \Es{mit bewusstem Verzicht auf zeit-
und personalaufwendige Maximalversorgung des Einzelnen zugunsten der
lebensrettenden Minimalversorgung Vieler.}

\subsubsection{Bei Gro{\ss}unfällen und Katastrophen versucht man das beste
herauszuholen}


\begin{tabularx}{\linewidth}{XX}
\hline
\TabHeading{Ziele}

\begin{itemize}
    \item das Bestmögliche
    \item für die grö{\ss}te Anzahl
    \item zur rechten Zeit
    \item am rechten Ort
\end{itemize}
 &
\TabHeading{Das ist nur erreichbar wenn:}

\begin{itemize}
    \item der Bestgeeignete auf jeder Stufe führt
    \item alle Einsatzkräfte Ihre Aufgabe kennen
    \item das dazu benötigte Material vorhanden, bereitgestellt und rasch
    einsatzbereit ist
    \item klare und eindeutige Führungsstrukturen vorhanden sind
    \item die Führungspyramide von allen strikt eingehalten wird
    \item die Organisation eingespielt/eingeübt ist
    \item Einsätze geplant und vorbereitet sind
\end{itemize}
\\\hline
\TabHeading{Dafür werden Sonderausrüstungen zusätzlich erforderlich,} da
ein erhöhter Bedarf besteht an:

\begin{itemize}
    \item Material
    \item Medikamenten
    \item Personal
\end{itemize}
&
\TabHeading{Wir müssen rechnen mit}
\begin{itemize}
    \item einem gro{\ss}en Patientenanfall
    \item besonderen Verletzungsmustern
    \item einer längeren Einsatzzeit vor Ort
    \item Überforderung der Helfer
\end{itemize}
\\\hline
\end{tabularx}

\subsubsection{Es gibt verschiedene Katastrophenarten}

\begin{center}
\begin{tabularx}{\linewidth}{XX}
\T{Naturkatastrophen}
 &
Erdbeben,\"Uberschwemmungen, Erdrutsche/Lawinen etc.
\\
\T{Technische Katastrophen}
 &
Flugzeugabsturz, Eisenbahnunglück, Reaktor\-unfall, Schiffs\-untergang
etc.
\\
\T{Sekundärkatastrophen}
 &
Hungersnot, Seuchen etc.
\\
\end{tabularx}
\end{center}

\subsubsection{Katastrophenhilfegesetze regeln die Zuständigkeiten}

Die Katastrophenhilfsgesetze regeln die Aufgaben der Länder,
Organisationen und Zuständigkeiten in Ihrem Bereich.
\Es{Ausschlie{\ss}lich eine Behörde erklärt ein au{\ss}ergewöhnliches
Schadensereignis zur Katastrophe} im Sinne des Gesetzes. Behördliche
Alarmpläne dienen zur Festlegung der Ma{\ss}nahmen der
Katastrophenhilfe und zur Koordination der im Anlassfall tätig
werdenden Einrichtungen. Je nach Ausdehnung des Schadensgebietes gibt
es unterschiedliche Kompetenzen:

\begin{center}
\begin{tabular}{m{2.375cm} c m{6.1590004cm}}
\toprule
\E{Gemeinde} 
 &
 ${\rightarrow}$
 &
Bürgermeister
\\
\E{Bezirk} 
 &
 ${\rightarrow}$
 &
Bezirkshauptmann
\\
\E{Bundesland}
 &
 ${\rightarrow}$
 &
Landeshauptmann
\\
\E{Bundesgebiet}
 &
 ${\rightarrow}$
 &
Innenminister%, Bundeskanzler 
\\\bottomrule
\end{tabular}
\end{center}

\subsubsection{Die Katastrophenwarnung mittels Sirene alarmiert die
Bevölkerung}

\includegraphics[width=\linewidth]{./images/noauth/khd-sirenen.png}

\subsection{Katastrophenschutz in Wien}
\begin{center}
\begin{tabular}{p{4.3190002cm}m{8.981cm}}
Wien ist für den Gro{\ss}\-schadens- und Katastrophenfall gut
gerüstet. Tausende Helfer sind dazu im Katastrophen\-schutz-Kreis
(\emph{{\quotedblbase}K-Kreis{\textquotedblleft}}) org\-anisiert.

Der K-Kreis ist ein Zu\-sammen\-schluss der beruflichen und freiwilligen
Hilfs- und Ein\-satz\-organisationen sowie wichtiger Behörden.
Dreh\-scheibe und gemeinsame Anlaufstelle ist dabei die Stadt Wien,
Magistrats\-direktion - Katastrophen\-manage\-ment und
Sofort\-ma{\ss}\-nahmen und die Helfer Wiens.
 &
\includegraphics[width=9.046cm,height=13.623cm]{./images/noauth/k-kreis08.png}
\CaptionofDiff{figure}{Der K-Kreis.
\SourcePicture{http://www.diehelferwiens.at/k-kreis , Copyright: Die Helfer
Wiens, Grafik: Erdle/Zimmermann/Schimanek}}{Der K-Kreis}
\\
\end{tabular}
\end{center}

\subsubsection{Katastrophenhilfe}

Hilfsma{\ss}nahmen die von uns als Sanitätsorganisation erwartet
werden:

\begin{itemize}
\item Bergung, Versorgung und Transport 
\item Pflegerische Betreuung 
\item Sorge für Unterkunft und Verpflegung 
\item Karitative Betreuung von Hilfsbedürftigen
\end{itemize}

\begin{center}
\begin{tabularx}{\linewidth}{lX}
\toprule
\TabHeading{Sanitätseinsatz}
 &
\begin{itemize}
     \item Retten/Bergen
     \item Triage
     \item Lebensrettende Sofortma{\ss}nahmen
     \item Erstversorgung
     \item Sanitätshilfe
     \item Abtransport
\end{itemize}
\\\midrule
\TabHeading{Pflegeeinsatz}
&
 Wenn die pflegerische Versorgung durch die bestehenden Einrichtungen
 nicht mehr sichergestellt ist.
 \begin{itemize}
     \item Erweiterung der vorhandenen Behandlungsmöglichkeiten
     \item Errichtung von Hilfskrankenhäusern
 \end{itemize}
\\\midrule
\TabHeading{Sozialeinsatz}
&
Wenn Menschen infolge einer Katastrophe ihre
Selbstversorgungsmöglichkeit verloren haben.

\begin{itemize}
    \item Einrichtung/Betrieb von Notunterkünften/Lagern
    \item Beschaffung, Zubereitung und Ausgabe von Lebensmitteln, Kleidern,
    Gebrauchsgegenständen
    \item Fürsorgliche Hilfe, Registrierung und Betreuung
\end{itemize}
\\\bottomrule
\end{tabularx}
\end{center}


\subsection{Grundsätze des Einsatzes}

\Fazit{Grundkomponenten: Führung -- Zeit -- Raum} 

Ziel ist es ohne überstürztes Handeln die Ausdehnung der Krisenlage
auf benachbarte und andere Räume zu verhindern. Dabei wird
gemä{\ss} der gesetzlichen Bestimmungen von Kräften der
Einsatzorganisationen, Behörden und anderen Einrichtungen koordiniert
vorgegangen.

\paragraph{Führung}
\Definition{Steuerndes Einwirken auf andere Menschen um ein bestimmtes Ziel zu
erreichen. }

Die Führung ist ein unentbehrliches Steuerungsinstrument zur
Erfüllung des Einsatzauftrages
({\quotedblbase}Ziel{\textquotedblleft}) und dem richtigen Einsatz von
Personal, Material und Transportmitteln 

Im Schadensraum gilt: {\quotedblbase}\Es{Ein einziger
Verantwortlicher}{\textquotedblleft} mit Entscheidungsbefugnis und den
nötigen Führungsmittel \E{({\quotedblbase}Einheit der
Führung{\textquotedblleft})}. Dieser wirkt mittels möglichst
einfacher Ma{\ss}nahmen \E{({\quotedblbase}Einfachheit{\textquotedblleft})}
auf den Einsatzverlauf ein
\E{({\quotedblbase}Handlungsfreiheit{\textquotedblleft})}. Mittel und
Mannschaften werden dabei gemä{\ss} ihrer Eigenart und
Leistungsfähigkeit eingesetzt \E{({\quotedblbase}\"Okonomie der
Kräfte{\textquotedblleft})}.Er muss rasch auf \"Anderungen der
Situation reagieren \E{({\quotedblbase}Beweglichkeit{\textquotedblleft})}
und gleichzeitig vorausschauend handeln
\E{({\quotedblbase}Reservenbildung{\textquotedblleft})}.

\Cave{{\quotedblbase}Führen{\textquotedblleft} und
{\quotedblbase}Retten{\textquotedblleft} ist nicht gleichzeitig
möglich!}

\paragraph{Sanitäter und Notärzte müssen Führungsaufgaben
übernehmen}~

Der erste am Schadensort eintreffende Sanitäter muss:

\begin{enumerate}
    \item Lagemeldung abgeben
    \item den Standort der örtlichen Einsatzleitung festlegen
    \item mit bereits anwesenden Einsatzkräften die ersten Ma{\ss}nahmen
    zur Rettung festlegen
    \item festlegen, wohin die Patienten im Schadensraum gebracht werden 
    sollen (Beginn einer Triage)
\end{enumerate}

\paragraph{Zeit}

Jedes Großschadens- oder Katastrophenereignis zeigt unabhängig von seiner
Ursache i.d.R. folgenden Zeitablauf:

\begin{table}[H]
\Captionof{table}{Phasen des Einsatzes.}

\begin{tabularx}{\linewidth}{XXX}
\TabHeading{1. Phase}
 &
\TabHeading{2. Phase}
 &
\TabHeading{3. Phase} 
\\\midrule\addlinespace
Isolation (Minuten)
 &
Retten (Stunden)
 &
Wiederherstellung (Tage)
\\\addlinespace
Schaden begrenzen  

- Lage erkennen 

- Melden
 &
Ordnung schaffen 

- Kommandoposten
 &
Lokale Einsatzleitung 

- Verbindung 

- Lage beurteilen
\\\addlinespace
\"Uberleben
 &
Weiterleben
 &
Endbehandlung
\\\addlinespace
Laien: 

- Nothilfe/Selbsthilfe 

- Retten / Erste Hilfe
 &
Geschulte Helfer: 

- LRSM 

- Transportfähigkeit
 &
Spezialisten: 

- Spital nach Dringlichkeit behandeln
\\\addlinespace
\multicolumn{3}{c}{\includegraphics[width=0.5\linewidth,height=1em]{./images/khd/arrow-right.png}}
\\\addlinespace\bottomrule
\end{tabularx}
\end{table}


\section[Sanitätsdienstliche Organisation im
Schadensraum und in der SanHist]{Hilfe braucht Raum: Sanitätsdienstliche Organisation im
Schadensraum und in der Sanitätshilfsstelle (SanHist)}
% TODO : Layout-Anpassung
\begin{anfxwarning}[author=GABS]{Layout-Anpassung notwendig.}
   Dieser Teil entspricht nicht dem Standard-Layout!
\end{anfxwarning}
Der Einsatz allgemein wird unterteilt in: 

\begin{itemize}
\item Schadensraum
\item Transportraum
\item Hospitalisationsraum
\end{itemize}

\begin{description}
    \item[Der Hospitalisationsraum] betrifft die Transportziele und
Aufnahmespitäler. 
\item[Der Transportraum] umfasst in diesem Zusammenhang alle Mittel und Wege um
die Patienten vom Schadens\-raum zum Hospitalisationsraum zu
verbringen. Dieser {\quotedblbase}allgemeine
Transportraum{\textquotedblleft} ist nicht zu ver\-wechseln mit dem
Transportraum der Sanitätshilfestelle (s.u.). 
\item[Der Schadensraum] ist der Ort des eigentlichen (schadens-)Geschehens. Es
gibt in ihm eine einheitliche Organisations\-form für den
Sanitätsdienst: Die \T{Sanitätshilfsstelle} (\T{SanHiSt}), welche aus den
3 Räumen \T{Triageraum}, \T{Behandlungsraum} und \T{Transportraum} besteht.
Definiert wird Sie in der ASB\"O Rahmen\-vor\-schrift für
Gro{\ss}unfälle und den Durchführungsbestimmungen für
Gro{\ss}unfälle. 
\end{description}



 



\begin{minipage}{12.605cm}
\includegraphics[width=\linewidth]{./images/khd/sanhist-plan-edited.png}
\Captionof{figure}{Gliederung des Einsatzraumes}
\begin{multicols}{2}

\begin{description}
    \item[\T{\"Aussere Absperrung}:] Gro{\ss}räumige Verkehrsumleitung
    \item[\T{Innere Absperrung}:] hält Schaulustige fern, hält Personen in Panik
    zurück
    \item[\T{Sicherheitsring}:] Absperrung des Schadenplatzes, Sicherheit für
    Einsatzkräfte
    \item[\T{Pforte}:] Einbahnsystem für Einsatz-Kfz, keine Fremdfahrzeuge,
    möglichst nur eine Pforte
    \item[\T{Leiter Information}] (i) für Weiterleitung von Informationen an die
    Au{\ss}enwelt (Presse, Angehörige, ...)
    \item[\T{Bergetriage}] (BT)
    \item[\T{Sanitätshilfsstelle} (\T{SanHist}):] mit Triageraum, Behandlungsraum und
    Transportraum;
    \item[\T{Einsatzleiter}] (EL)
    \item[\T{Leitender Notarzt}] (LNA)
    \item[\T{Mobile Leitstelle}] (MLS)
    \item[\T{Material- und Meldestelle}] (M)
    \item[\T{Sammelstelle Tote}] (t) und
    \item[\T{Sammelstelle Unverletzte}] (U)
\end{description}


\end{multicols}
\end{minipage}

\subsection{Triageraum}

\Z{\emph{Triage (frz.) = Sichtung, Auswahl}}

\Definition{Triage ist das Begutachten aller Patienten mit geringen Mitteln und
Aufwand sowie das Einteilen in Triagegruppen mit dem Ziel möglichst
viele Leben zu retten!}

{\quotedblbase}Life before Limbs{\textquotedblleft} = Leben vor
Gliedma{\ss}en, Funktion vor Schönheit!

\subsubsection{Bergetriage}~


\Parallel{
Nur wenn keine Gefahrenzone und genügend Personal/Zeit vorhanden ist!

\begin{itemize}
    \item BT Team muss mind. aus einem NFS bestehen
    \item Gefahren beachten
    \item Erst Sichtung der Betroffenen und Festlegung der Dringlichkeit
    \item Nicht behandeln!
\end{itemize}
}{
\begin{center}
\includegraphics[width=2cm]{./images/noauth/sanhist-bergetriage.png}
\end{center}
}


\subsubsection{Triage (Sichtung)}~

%%%%%%%%%%%%%%%%%%%%%%%%%%%%%%%%%%%%%%%%%%%%%%%%%%%%%%%%%%%
% TODO Dynamik I-IV: Überarbeitung notwendig.
\begin{anfxwarning}[author=GABS]{Dynamik I-IV: Überarbeitung notwendig.}
   Dynamik I-IV: Überarbeitung notwendig. Aussage der i-IV-Dynamik kommt nicht
   gut herus.
\end{anfxwarning}
%%%%%%%%%%%%%%%%%%%%%%%%%%%%%%%%%%%%%%%%%%%%%%%%%%%%%%%%%%%


\Parallel{
\begin{itemize}
        \item Zentrale Registrierung aller Patienten mittels Patientenleitsystem
    (PLS)
        \item Zuordnung der Patienten zu Triagegruppen
    \begin{description}
        \item[\Code{I~~}] Sofortbehandlung, lebensrettende Noteingriffe
        \item[\Code{II~}] Dringende Behandlung, Schwerverletzte
        \item[\Code{III}] Minimalbehandlung
        \item[\Code{IV~}] Abwartende Behandlung
    \end{description}
    \item Bei Missverhältnis zwischen der Zahl der Betroffenen und den
    Be\-handlungs\-möglich\-keiten
    \item Schnell, mit geringem Aufwand, hoher Zeitdruck

  Maximaler Zeitaufwand pro Patient gesamt:
    \begin{itemize}
        \item Gehender 1 Minute
        \item Liegender 3 Minuten
    \end{itemize}    
    \item Keine detaillierte Diagnostik möglich
\end{itemize}
}{
\begin{center}
% \includegraphics[width=2.cm]{./images/noauth/sanhist-triage.png}
\TODO{GABS}{2010-04-27}{Foto Triageplatz}{}
\end{center}
}

Die Zuordnung zu einer Triagegruppe kann sich im Verlauf verändern,
    eine ständige Nachtriage ist erforderlich! Änderungen ergeben sich besonders
    durch:

\begin{enumerate}
    \item \Margin{z.\,B. III → I}Verschlechterung des Patientenzustandes.
    \item \Margin{z.\,B. I → III}Verbesserung des Patientenzustandes (Behandlung).
    \item \Margin{z.\,B. IV → I}Verbesserung der Versorgung wegen nachgerücktem 
    Personal und Material.
\end{enumerate}

\Parallel
{Zwei Triageplätze, damit der NA abwechselnd
triagieren kann
\begin{itemize}
    \item Vorbereitung durch RS/NFS
    \item Entkleiden
    \item Vitalparameter erheben %\footnote{RR, HF, AF, Sp$O_2$}
    \item Untersuchung durch den Notarzt
\end{itemize}}
{
\includegraphics[width=5.132cm,height=3.783cm]{./images/noauth/sanhist-tragestelle-pic.jpg}
}




\subsection{Behandlungsraum}



\begin{tabulary}{\linewidth}{lL}
\toprule
\begin{minipage}{3.5cm}
\includegraphics[width=1.489cm,height=1.61cm]{./images/noauth/sanhist-behandlung.png}
\end{minipage}
 &
Es werden drei Behandlungsstellen eingerichtet
\\
\begin{minipage}{3.5cm}
\includegraphics[width=1.629cm,height=1.489cm]{./images/noauth/sanhist-triage-1}
\end{minipage}


 &
Sofortbehandlung vor Ort

T1 Lebensrettende Soforteingriffe 
\\
\begin{minipage}{3.5cm}
\includegraphics[width=1.669cm,height=1.578cm]{./images/noauth/sanhist-triage-2.png}
\end{minipage}
 &
Dringende Behandlung

T2 Schwerverletzte
\\
\begin{minipage}{3.5cm}
\includegraphics[width=1.601cm,height=1.587cm]{./images/noauth/sanhist-triage-3.png}
\includegraphics[width=1.587cm,height=1.641cm]{./images/noauth/sanhist-triage-4.png}
\end{minipage}

 &
Warten

T3 Warten Leichtverletzte

T4 Warten Schwerverletzte
\\\bottomrule
\end{tabulary}




\subsubsection{I -- Sofortbehandlung vor Ort}

Alle Patienten, welche ohne sofortige Behandlung sterben bzw. schwere
Schäden erleiden

\begin{flushleft}
\begin{tabular}{m{5.642cm}m{4.379cm}m{3.076cm}}
Welche Patienten?

\begin{itemize}
    \item Atemstörungen
    \item Spannungspneumothorax
    \item schwere äu{\ss}ere Blutungen
    \item schwerer traumatischer Schock
\end{itemize}
&
Ma{\ss}nahmen

\begin{itemize}
    \item \"Uberprüfung der Triage
    \item Volumenersatz
    \item Intubation
    \item Thoraxdrainage
    \item Noteingriffe
\end{itemize}
&
\includegraphics[width=3cm]{./images/noauth/sanhist-triage-1.png}
\\
\end{tabular}
\end{flushleft}


\subsubsection{II -- Dringende Behandlung}

Behandlung und Herstellung der Transportfähigkeit von schwer
verletzten/erkrankten Personen, eine weitere Behandlung im Krankenhaus
ist i.d.R. notwendig.

\begin{flushleft}
\begin{tabular}{m{5.642cm}m{4.379cm}m{3.076cm}}
Welche Patienten?

keine unmittelbare Lebensbedrohung

\begin{itemize}
    \item SHT
    \item Wirbelverletzungen mit Lähmungen
    \item Bauchtrauma, innere Blutungen
    \item schwere Thoraxtraumen
    \item gro{\ss}e Gefä{\ss}verletzungen
    \item offene Knochenbrüche und  Gelenks\-verletzungen
    \item Augenverletzungen
    \item gro{\ss}e Weichteilverletzungen
    \item Verbrennungen 15 - 30\%
    \item geschlossene Knochenbrüche und Gelenks\-ver\-letzungen
\end{itemize}
&
Ma{\ss}nahmen

\begin{itemize}
    \item Verbände
    \item Schockbehandlung
    \item Fixierung/Schienung
    \item Schmerzbekämpfung
    \item pflegerische Ma{\ss}nahmen
    \item LRSM wenn nötig
\end{itemize}
&
\includegraphics[width=3cm]{./images/noauth/sanhist-triage-2.png}

\\
\end{tabular}
\end{flushleft}


\subsubsection{III -- Warten Leichtverletzte}

Alle Patienten, welche nur einer Minimalbehandlung bedürfen.

\begin{flushleft}
\begin{tabular}{m{5.642cm}m{4.379cm}m{3.076cm}}
Welche Patienten?

\begin{itemize}
    \item Verbrennungen bis 15\%
    \item kleine Weichteilwunden
    \item einfache Knochenbrüche
    \item Prellungen, Zerrungen
    \item leichtes SHT
\end{itemize}
&
Ma{\ss}nahmen

\begin{itemize}
    \item \"Uberprüfung der Triage
    \item pflegerische Ma{\ss}nahmen
    \item Betreuung (Kriseninterventionsteam (\Abk{KIT})
    \item Entlassung ambulant Behandelter zur SaStU
\end{itemize}
&
\includegraphics[width=3cm]{./images/noauth/sanhist-triage-3.png}
\\
\end{tabular}
\end{flushleft}


\subsubsection{IV -- Warten Schwerverletzte}

Alle Patienten, welche mit dem zur Verfügung stehenden Mitteln
(Personal, Material) noch nicht versorgt werden können . Abwartende
Behandlung.

\begin{center}
\begin{tabular}{m{5.702cm}m{4.0550003cm}m{3.3439999cm}}
Welche Patienten?

\begin{itemize}
    \item schweres Polytrauma
    \item schwerstes \Abk{SHT}
    \item offene Körperhöhlentraumen
    \item Mehrhöhlenverletzungen
    \item Verbrennungen über 50\%
\end{itemize}
&
Ma{\ss}nahmen

\begin{itemize}
    \item \"Uberprüfung der Triage
    \item Behandlung  (Schmerzbekämpfung)
    \item Psychologische Betreuung, evt. durch \Abk{KIT}
    \item \"Uberwachung
\end{itemize}
 &
\includegraphics[width=3cm]{./images/noauth/sanhist-triage-4.png}
\\
\end{tabular}
\end{center}


\subsubsection{Die Transportpriorität ist unabhängig von der Triagegruppe}

Nach einer erfolgten Behandlung kann eine Transporpriorität zugewiesen
werden. Man unterscheidet hierbei: 

\begin{itemize}
\item A: Hohe Transportpriorität: Der rasche Transport zur
frühzeitigen Fachbehandlung nach ärztlicher Notversorgung ist
angesagt. 
\item B: Niedrige Transportpriorität: Der Transport zur Fachbehandlung
kann verzögert erfolgen. 
\end{itemize}






\subsection{Transportraum}

\subsubsection{Verladestelle}

\Parallel{
Sollte in der Nähe der Behandlungsstelle II errichtet werden.

Es muss protokolliert werden, welcher Patient mit welchem
Transportmittel in welches Krankenhaus eingeliefert wird.

Einbahnregelung!
}{
\begin{center}
\includegraphics[width=3cm]{./images/noauth/sanhist-verladestelle.png}
\end{center}
}



\subsubsection{Wagenhalte\-platz}


\Parallel{
\begin{itemize}
    \item Groß genug
    \item Schrägparkordnung, freies Heck zum Ent-/Beladen
    \item Sortiert nach Einsatzmittel (\Abk{NAW}, \Abk{RTW}, \ldots)
    \item Schlüssel stecken lassen
\end{itemize}
}{
\begin{center}
\includegraphics[width=3cm]{./images/noauth/sanhist-wagenhalteplatz.png}
\end{center}
}


Der Wagenhalteplatz muss \Es{frühzeitig} bestimmt werden, da weitere Kräfte
rasch nachfolgen werden. Er sollte in ausreichender \Es{Entfernung} festgelegt
werden um den späteren Aufbau der weiteren SanHiSt nicht zu behindern. Bedenke:
Triage- und Behandlungsraum (und die Verladestelle) brauchen viel Platz!





%  Falscher Wagenhalteplatz
% 
% 
% Warning: Image ignored
% Unhandled or unsupported graphics:
%\includegraphics[width=3.375cm,height=2.438cm]{AASdev20090228-img129}




\subsubsection{Hubschrauber\-landestelle}

\Parallel{
\begin{itemize}
    \item Die allgemeinen Vorschriften für Hubschrauberlandeplätze sind zu beachten
    \item Freie ebene Fläche
    \item Sichere Umgebung (Stromleitungen!)
    \item Weit genug weg vom Behandlungsraum (Luftwirbel!)
\end{itemize}
}{
~
% TODO BILD: Hubschrauberlandestelle
}






\subsubsection{Exkurs: Transportmanagement}~

\Parallel{
\begin{itemize}
    \item Spiralförmige Verteilung
    \item Zeitliche Staffelung
    \item Fachliche Zuordnung
    \item Katastrophe nicht ins Spital verlagern!
\end{itemize}
}{
\includegraphics[width=\linewidth]{./images/noauth/khd-transportmanagement.png}
}





\subsection{Sammelstellen}

\subsubsection{Sammelstelle Unverletzte}


\Parallel{
Ziel: Betreuung Unverletzter und Betroffener

\begin{itemize}
    \item Betreuen
    \item Beruhigen
    \item Informieren
    \item \Abk{KIT}\footnote{In Wien ist die Akutbetreuung Wien (\Abk{ABW})
    zuständig.} integrieren
\end{itemize}
}{
\begin{center}
\includegraphics[origin=t,width=2.01cm,height=2.154cm]{./images/noauth/sanhist-sammelstelle-unverletzte.png}
\end{center}
}


\subsubsection{Sammelstelle Tote}

\Parallel{
Sollte diskret eingerichtet werden.

Abgeschirmt von der Sammel\-stelle Unverletzte, dem Be\-handlungs\-raum
und der Au{\ss}en\-welt
}{
\begin{center}

\includegraphics[width=1.906cm,height=1.9cm]{./images/noauth/sanhist-sammelstelle-tote.png}
\end{center}
}





\subsection{Organisation}

\subsubsection{Material- / Meldestelle}

\Parallel{
\begin{itemize}
    \item Dient zur Registrierung des Personals und Materials
    \item Eintreffende Einheiten haben sich bei dieser zu melden
    \item Steht im engen Kontakt zum Leiter SanHist und EL
\end{itemize}
}{
\begin{center}
\includegraphics[width=2.176cm,height=2.126cm]{./images/noauth/sanhist-material-meldestelle.png}
\end{center}
}



\subsubsection{Mobile Leitstelle (MLS)}

\Parallel{
\begin{itemize}
    \item Die \Abk{MLS} beherbergt i.d.R. die Ein\-satz\-leitung sowie wichtige
    Administrations-, Tele\-kom\-munikations- und Funkeinrichtungen.
    \item Bei Bedarf werden u.U. hier Handfunkgeräte an die Bereichsleiter
    ausgegeben.
    \item MLS ist die Brücke zwischen verschiedenen Funkkanälen und
    zuständig.
    \item Wenn die MLS die Einsatzleitung beinhaltet wird sie mit einem
    rotem Dreh- oder Blinklicht markiert.
\end{itemize}
}{
\begin{center}
\begin{center}
\includegraphics[width=\linewidth]{./images/khd/khd-mls.jpg}
\end{center}
\end{center}
}


\subsection{Führungsstruktur in der SanHist}

\begin{center}
\begin{tabular}{m{1.8cm}m{4.5cm}|m{1.8cm}m{4.5cm}}
\includegraphics[width=1.654cm,height=1.624cm]{./images/noauth/sanhist-einsatzleiter.png}
&
\Es{Einsatzleiter} (Rettungsdienst) 

Er ist der Vorgesetzte für die
eingesetzten Sanitäter und sonstigen sanitätsdienstlichen Hilfskräfte. Er hat
\E{keine} direkte Weisungsbefugnis gegenüber anderen eingesetzten Organisationen. Er kann Ärzten
\E{keine Weisungen in medizinischen Fragen} geben. 
&
% Unhandled or unsupported graphics:
\includegraphics[width=1.62cm,height=1.62cm]{./images/noauth/sanhist-material-meldestelle.png}
&
\Es{Leiter Material-/Meldestelle}
\\ 
% Unhandled or unsupported graphics:
\includegraphics[width=1.62cm,height=1.62cm]{./images/noauth/sanhist-leiter-sanhist.png}
&
\Es{Leiter Sanhist}
&
% Unhandled or unsupported graphics:
\includegraphics[width=1.489cm,height=1.61cm]{./images/noauth/sanhist-behandlung.png}
&
\Es{Leiter Behandlung}
\\ % Unhandled or unsupported graphics:
\includegraphics[width=1.62cm,height=1.62cm]{./images/noauth/sanhist-lna.png}
&
\Es{LNA Leitender Notarzt} 

Der LNA ist der Vorgesetzte der eingesetzten Ärzte.
& 
% Unhandled or unsupported graphics:
\includegraphics[width=1.62cm,height=1.62cm]{./images/noauth/sanhist-leiter-transport.png}
&
\Es{Leiter Transport}
\\ 
% Unhandled or unsupported graphics:
\includegraphics[width=1.62cm,height=1.62cm]{./images/noauth/sanhist-leiter-triage.png}
&
\Es{Leiter Triageraum/Triage}
& 
% Unhandled or unsupported graphics:
\includegraphics[width=1.62cm,height=1.62cm]{./images/noauth/sanhist-information.png}
&
\Es{Leiter Information}
\\
\end{tabular}
\end{center}


\section{Patientenleitsystem (PLS)}
% TODO : Layout-Anpassung
\begin{anfxwarning}[author=GABS]{Layout-Anpassung notwendig.}
   Dieser Teil entspricht nicht dem Standard-Layout!
\end{anfxwarning}

Das \index{Patientenleitsystem}\T{Patientenleitsystem} (\index{PLS}\T{PLS})
ermöglicht eine \Es{eindeutige Kennzeichnung} von Patienten und Zuordnung
von Behandlungsma{\ss}nahmen und Befunden. Der wichtigste Bestandteil
ist die \index{Patientenleittasche}\T{Patientenleittasche}
(\index{PLT}\T{PLT}). 

Sie besteht aus einer orangenen Hülle welche \E{fortlaufend nummeriert}
ist. Auf dieser Kunststoffhülle können wichtige Notizen wie
Triagegruppe, etc., vorgenommen werden. In der Hülle befinden sich
das \E{Behandlungsprotokoll} (blau), \E{Identifikationsprotokoll} (rosa),
\E{Zusatzkarten}, \E{Kontaminationsaufkleber} und ein Bogen mit gleichlautend
\E{nummerierten Klebeetiketten}. Am unteren Ende befinden sich \E{zwei
Abschnitte} welche die Rückverfolgbarkeit des Transports
sicherstellen. Die PLT ist mit einer Gummischnur versehen damit sie dem
Patienten umgehängt werden kann. 

\Es{\E{Spätestens} an der Triagestelle} bekommt jeder Patient oder
Betroffene {\quotedblbase}seine{\textquotedblleft} \Abk{PLT}. Auch Unverletzte
und Verstorbene werden mittels \Abk{PLS} registriert.

\paragraph{Zusatzkarten}

Bei einer eventuellen Bergetriage bekommt jeder Patient eine
Patientenleittasche (\Abk{PLT}). Patienten, welche \Es{bevorrangt gerettet}
werden sollen, werden zusätzlich mit der gelben
{\fcolorbox{Black}{Yellow}{\sffamily\bfseries DRINGEND}}\,-Karte
gekennzeichnet. Die Karte wird auf der Triagestelle wieder entfernt. Tote
erhalten zusätzlich zur \Abk{PLT} eine schwarze
{\fcolorbox{Black}{White}{\sffamily\bfseries VERSTORBEN}}\,-Karte.

\paragraph{Behandlungsprotokoll / Identifikationsprotokoll}

Die Protokolle werden im Behandlungsraum ausgefüllt wenn Zeit dafür
ist. Das Ausfüllen der Protokolle darf den Abtransport unter keinen
Umständen verzögern. Die Protokolle können auch leer bleiben.

\paragraph{Kontaminationsaufkleber}\index{Kontaminationsaufkleber}

Wenn der Patient kontaminiert ist, wird ein Kontaminationsaufkleber in
Form eines gelben Dreiecks mit schwarzem Rand auf die \Abk{PLT} geklebt. Dies soll
alle nachfolgenden Personen, welche mit dem Patienten in
Kontakt treten, warnen. Die Aufkleber sind auf \Es{beiden
Seiten} der \Abk{PLT} anzubringen!

\paragraph{PLT-Nummer}

Die \Abk{PLT} und damit der Patient ist durch eine \Es{eindeutige, einmalige
Nummer} \Es{identifiziert}. Mittels \Es{Selbstklebeetiketten} können die
Nummern der Patienten in verschiedenen Listen schnell erfasst werden. Das ist
in vielen Bereichen besonders wichtig (Triage, Behandlung, Transport). Außerdem
können alle Effekte des Patienten damit gekennzeichnet werden.

Die Patientenleittasche verbleibt solange am Patienten, \E{bis das
Aufnahmeverfahren im Krankenhaus abgeschlossen ist}, d.h. bis der Patient im
Spital administrativ erfasst ist. Anschlie{\ss}end wird die \Abk{PLT} der
Krankengeschichte des Patienten beigefügt.

\Fazit{Die Nummerierung der \Abk{PLT} ermöglicht die Zuordnung und Erfassung von
Patienten. Das ist eine der Hauptaufgaben der \Abk{PLT}!}

\A{PLT-Fotos \ldots}


\includegraphics[width=\linewidth]{./images/khd/plt-inhalt-beschriftet.jpg}

\begin{multicols}{2}

\begin{description}
    \item[a] Zusatzkarte {\fcolorbox{Black}{White}{\sffamily\bfseries
VERSTORBEN}}
    \item[b] Zusatzkarte 
{\fcolorbox{Black}{Yellow}{\sffamily\bfseries DRINGEND}}
    \item[c] Klebeetiketten mit PLT-Nummer
    \item[d] Aufkleber für kontaminierte
Gegenstände 
    \item[e] Behandlungsprotokoll (blau)
    \item[f] Identifikationsprotokoll (rosa)
\end{description}

\end{multicols}

 
\subsubsection{PLT Vorderseite}

\Parallel{
\paragraph{Feld \emph{"`Triage"'}} \textcircled{a}

Kennzeichnung der Triagegruppe:

Bei Triage nur Einteilung in {\sffamily I}, {\sffamily II},
{\sffamily III} oder {\sffamily IV}

\paragraph{Zweiter Triageblock} \textcircled{b}

Dient zur 
\begin{itemize}
    \item Umtriage oder
    \item Spitalstriage
\end{itemize}

\paragraph{Feld \emph{"`Transport"'}} \textcircled{c}

Kennzeichnung der Transportpriorität, wird später
entschieden

\begin{description}
    \item[{\sffamily A}: Hohe Transportpriorität:] Der
    rasche Transport zur frühzeitigen Fachbehandlung nach
    ärztlicher Notversorgung ist angesagt.
    \item[{\sffamily B}: Niedrige Transportpriorität:] Der Transport
zur Fachbehandlung kann verzögert erfolgen. 
\end{description}

  

 




}{
\includegraphics[width=\linewidth]{./images/khd/plt-vorderseite-beschriftet-50pc.jpg}
}

%%%%%%%%%%%%


\begin{minipage}[t]{0.48\linewidth}
\includegraphics[width=\linewidth]{./images/khd/plt-vorderseite-umtriagierung15-50pc.jpg}
\Captionof{figure}{PLT: Umtriagierung innerhalb der ersten
Viertelstunde}
\begin{itemize}
\item Mittels Pfeil von einem zum nun richtigen Triagefeld
\item In der gleichen Zeile
\end{itemize}
\end{minipage}
\hspace{0.04\linewidth}
\begin{minipage}[t]{0.48\linewidth}
\includegraphics[width=\linewidth]{./images/khd/plt-vorderseite-umtriagierung60-50pc.jpg}
\Captionof{figure}{PLT: Umtriagierung nach längerer Zeit}
\begin{itemize}
    \item Neues Kreuz beim richtigen Triagefeld
    \item in der unteren Zeile
\end{itemize}
\end{minipage}
	





\Parallel{
\paragraph{Rückseite}~

Der Triagearzt trägt auf der Rückseite die zu setzenden
\Es{Ma{\ss}nahmen} ein:

\begin{itemize}
    \item Atmung ($O_2$, Intubation, usw..)
    \item Kreislauf (Blutstillung usw...)
    \item Medikamente
    \item Lagerung
\end{itemize}

Der Durchführende (Arzt, NFS, RS) bestätigt die Durchführung der
Ma{\ss}nahmen im rechten Feld durch eine Paraphe.

\vspace{10em}

\paragraph{Abriss für das Zielspital:}~

Wird bei der Einlieferung mit der Ankunftszeit versehen und verbleibt in
der Krankengeschichte.

Abriss für die SanHist -- Leiter Transport:

Eintragen: Zielspital, Kfz Nummer, Verladezeit

Verbleibt beim Leiter Transport/Protokollführer
}{
\includegraphics[width=\linewidth]{./images/khd/plt-rueckseite-angekreuzt-50pc.jpg}
}










\clearpage
\section{Gefahrengutunfall}
\label{topic:gefahrengutunfall}
% TODO : Layout-Anpassung
\begin{anfxwarning}[author=GABS]{Layout-Anpassung notwendig.}
   Dieser Teil entspricht nicht dem Standard-Layout!
\end{anfxwarning}


\Parallel{
% \includegraphics[width=\linewidth]{./images/noauth/AASdev20090228-img149.jpg}
\includegraphics[width=\linewidth]{./icons/under-construction.png}
}{
% \includegraphics[width=\linewidth]{./images/noauth/AASdev20090228-img150.jpg}
\includegraphics[width=\linewidth]{./icons/under-construction.png}
}

\TODO{GABS}{2010-04-27}{Bilder 2x Gefahrengutunfall/LKW}{}


\Parallel{
\TabHeading{Mögliche Gefahren für den Sanitäter}

\begin{itemize}
    \item Brand bzw. Explosionsgefahr
    \item Entweichen unter Druck stehender Gase
    \item Giftigkeit, Ansteckungsgefahr
    \item \"Atzwirkung
    \item Radioaktivität
    \item Umweltgefährdung
\end{itemize}
}{
\TabHeading{Gefahrenquellen können entstehen bei}

\begin{itemize}
    \item Unfälle bei der Herstellung
    \item Unfälle bei der Verwendung von gefährlichen Stoffen
    \item Transportunfälle
\end{itemize}
}




\begin{tabulary}{\linewidth}{LL}\toprule
\multicolumn{2}{c}{\TabHeading{Andere Gefahrenquellen:}}
\\\midrule
Landwirtschaft
 &
$CO_2$ in Gärkellern, Siloanlagen, Brunnenschächte

Pflanzenschutzmittel
\\\addlinespace
Klärgruben
 &
$CO_2$, Methan, Schwefelwasserstoff
\\\addlinespace
Obstlagerhallen
 &
wenig $O_2$, viel $CO_2$ oder $N_2$ halten Obst frisch
\\\addlinespace
Brände in Wirtschaftsgebäuden
 &
extrem gefährliche Rauchgase
\\\addlinespace
Tiefgaragen und Tunnels
 &
erhöhte $CO$- und $CO_2$-Werte
\\\addlinespace
Flüssiggas
 &
Campingplätze, feste und mobile Küchen
\\\addlinespace
Pools, Schwimmbäder
 &
Chlor (Wasserdesinfektion)
\\\addlinespace
Industrie
 &
Kühlmittel in gro{\ss}en Kälteanlagen (Ammoniak)
\\\addlinespace
Speziallabors, Einrichtungen mit und zur Verarbeitung von infektiösen
Abfällen oder Tierkadaver
 &
Infektionsgefahr
\\\addlinespace
Nuklearmedizinische Institute, Industrie, Forschungseinrichtungen
 &
Radioaktivität
\\\bottomrule
\end{tabulary}


\subsection{Allgemeine Einsatzrichtlinie}
% TODO : Layout-Anpassung
\begin{anfxwarning}[author=GABS]{Layout-Anpassung notwendig.}
   Dieser Teil entspricht nicht dem Standard-Layout!
\end{anfxwarning}



\begin{description}
    \item[G]efahr erkennen
    \item[A]bsperrma{\ss}nahmen durchführen
    \item[S]pezialkräfte anfordern
\end{description}







%\begin{gWideParEnv}
%\begin{tabular}{p{0.2\linewidth}p{0.5\linewidth}p{0.3\linewidth}}
\begin{tabular}{p{3cm}p{\linewidth - 7.6cm}p{3.1cm}}
\toprule
\noindent\TabHeading{{\color{ubuntured}G}efahr erkennen}
 &
\noindent Die Suche nach Warntafeln ist bei umgestürzten Fahrzeugen erschwert
 &
~

\includegraphics[width=3cm]{./images/noauth/gefahrgut-unfall.jpg}
\\
\noindent\TabHeading{{\color{ubuntured}A}bsperr\-ma{\ss}nahmen
durch\-führen} &
Die Absperrung sollte gro{\ss}räumig  sein, genügender Platz für
die technischen Hilfeleistungen muss einkalkuliert werden. Beachte sich
ändernde Situationen  (Windrichtung, Niederschlag)
 &
~

\includegraphics[width=3cm]{./images/noauth/gefahrgut-absperrung.png}
\\
\noindent\TabHeading{{\color{ubuntured}S}pezial\-kräfte anfordern}
 &
Informationen angeben (Kemmler-Nummer, etc.)
 &
~

\includegraphics[width=3cm]{./images/noauth/gefahrgut-feuerwehr.jpg}
\\\bottomrule
\end{tabular}
%\end{gWideParEnv}



\subsection{Gefahr erkennen: Kennzeichnung von gefährlichen Stoffen und
Gefahrenbereichen}
% TODO : Layout-Anpassung
\begin{anfxwarning}[author=GABS]{Layout-Anpassung notwendig.}
   Dieser Teil entspricht nicht dem Standard-Layout!
\end{anfxwarning}

\begin{center}
\begin{tabulary}{\textwidth}{LLLL}
\toprule
\multicolumn{2}{c}{\TabHeading{Hinweis auf Gefahr}} 
&
\multicolumn{2}{c}{\TabHeading{Warnzeichen}}
\\\midrule
% Unhandled or unsupported graphics:
\includegraphics[width=1.475cm,height=1.452cm]{./images/noauth/gefahr-explosion.png}
 &
Explosionsgefährlich
 &
\includegraphics[width=1.541cm,height=1.426cm]{./images/noauth/warnung-allgemein.png}
 &
Allgemeine Gefahr
\\
\includegraphics[width=1.452cm,height=1.414cm]{./images/noauth/gefahr-giftig.png}
 &
Giftig
 &
\includegraphics[width=1.673cm,height=1.414cm]{./images/noauth/warnung-radioaktiv.png}
 &
Radioaktive Stoffe, ionisierende Strahlen
\\
\includegraphics[width=1.497cm,height=1.611cm]{./images/noauth/gefahr-hochentzuendlich.png}
 &
Hochentzündlich
 &
\includegraphics[width=1.696cm,height=1.412cm]{./images/noauth/warnung-gasflaschen.png}
 &
Gasflaschen
\\
\includegraphics[width=1.475cm,height=1.435cm]{./images/noauth/gefahr-reizend.png}
 &
Reizend
 &
\includegraphics[width=1.673cm,height=1.414cm]{./images/noauth/warnung-laser.png}
 &
Laserstrahlen
\\
\includegraphics[width=1.488cm,height=1.599cm]{./images/noauth/gefahr-aetzend.png}
 &
\"Atzend
 &
 &
\\\hline
\end{tabulary}
\end{center}

\subsubsection{Transport gefährlicher Güter}

Kennzeichnungspflicht besteht erst dann, wenn die transportierte Menge
eines Stoffes die gesetzlich erlaubte Mindestmenge überschreitet.
Diese Mindestmenge gilt für jeden einzelnen Stoff und nicht für die
Gesamtmenge.

Es können also erhebliche Mengen gefährlicher Güter ohne
Kennzeichnung des Fahrzeuges transportiert werden. 

\subsubsection{Gefahrzettel}

Bei Stückguttransporten wird jedes Versandstück gekennzeichnet. An
Fahrzeugen, die kein Stückgut transportieren, müssen Gefahrenzettel
an den Längsseiten und an der Rückseite angebracht werden.

\begin{flushleft}
\begin{tabular}{m{2.cm}m{3.5cm}m{2.cm}m{3.5cm}}
 Beispiele
 &
 &
 &
\\
\includegraphics[width=1.659cm,height=1.885cm]{./images/noauth/gefahrzettel-feuergefaehrlich-fluessig.png}
 &
Feuergefährlich \par(flüssige Stoffe)
 &
\includegraphics[width=1.856cm,height=2.096cm]{./images/noauth/gefahrzettel-selbstentzuendlich.png}
 &
Selbstentzündlich
\\
\includegraphics[width=1.835cm,height=1.902cm]{./images/noauth/gefahrzettel-feuergefaehrlich-fest.png}
 &
Feuergefährlich \par(feste Stoffe)
 &
\includegraphics[width=1.909cm,height=2.162cm]{./images/noauth/gefahrzettel-enzuendlich-wasser.png}
 &
Entzündbare Gase bei Berührung mit \Es{Wasser}
\\
\end{tabular}
\end{flushleft}

\subsubsection{Gefahrentafel}

\begin{table}[H]
\Captionof{table}{Kennzeichnung durch Gefahrentafel.}

\begin{tabular}{m{3.225cm}m{3.225cm}m{3.225cm}m{3.225cm}}
\multicolumn{2}{m{6.65cm}}{\centering Allgemeine Kennzeichnung (Sammeltransporte)} 
&
\multicolumn{2}{m{6.65cm}}{\centering Spezielle
Kennzeichnung} 
\\\midrule
\includegraphics[width=\linewidth]{./images/khd/warntafel-allgemein.png}
 &
 &
\includegraphics[width=\linewidth]{./images/khd/warntafel-un-60-1710.png}
 &
Gefahrnummer\par (\T{Kemler-Nummer})

Stoffnummer\par (\T{UN-Nummer})
\\\bottomrule
\end{tabular}
\end{table}


\subsubsection{Gefahrnummer}

\begin{table}[H]
\Captionof{table}{Gefahrennummer.}

\begin{tabular}{cp{11cm}}
\addlinespace
 2 & Entweichen von Gas durch Druck oder chemische Reaktion
\\\addlinespace
 3 & Entzündbarkeit von Flüssigkeiten (Dämpfen) und Gasen
\\\addlinespace
 4 & Entzündbarkeit fester Stoffe
\\\addlinespace
 5 & Oxidierende (brandfördernde) Wirkung
\\\addlinespace
 6 & Giftig
\\\addlinespace
 7 & Radioaktivität
\\\addlinespace
 8 & \"Atzwirkung
\\\addlinespace
 9 & Gefahr einer spontan heftigen Reaktion
\\\addlinespace\bottomrule
\end{tabular}
\end{table}


Die Verdopplung einer Ziffer weist auf die Zunahme der entsprechenden 
Gefahr hin.

X  vor der Nummer zur Kennzeichnung der Gefahr ${\rightarrow}$ Stoff
reagiert in  gefährlicher Weise mit Wasser

Wenn die Gefahr eines Stoffes ausreichend von einer einzigen Ziffer 
angegeben werden kann, wird dieser Ziffer eine Null angefügt.

Folgende Ziffernkombinationen haben eine besondere Bedeutung

\begin{table}[H]
\Captionof{table}{Spezielle Zifferkombinationen bei der Gefahrennummer.}

\begin{tabular}{cp{11cm}}
\addlinespace
22 & tiefgekühltes Gas
\\\addlinespace
X323 & entzündbarer flüssiger Stoff, der mit Wasser gefährlich reagiert, wobei entzündbare Gase entweichen
\\\addlinespace
X333 & selbstentzündliche Flüssigkeit, die mit Wasser gefährlich reagiert
\\\addlinespace
X423 & entzündbarer fester Stoff, der mit Wasser gefährlich reagiert, wobei brennbare Gase entweichen
\\\addlinespace
44 & entzündbarer fester Stoff, der sich bei erhöhter Temperatur in Geschmolzenem Zustand befindet
\\\addlinespace
539 & entzündbares organisches Peroxid
\\\addlinespace
90  & verschiedene gefährlich Stoffe
\\\addlinespace\bottomrule
\end{tabular}
\end{table}


\subsection{Absperrmaßnahmen durchführen}
% TODO : Layout-Anpassung
\begin{anfxwarning}[author=GABS]{Layout-Anpassung notwendig.}
   Dieser Teil entspricht nicht dem Standard-Layout!
\end{anfxwarning}

\TakeNotesLarge

\begin{table}[h]
\begin{gWideParEnv}
\Captionof{table}{Bilderserie \SlideHeading{Absperrung}.}

\begin{tabularx}{\linewidth}[H]{XXX}
\includegraphics[width=\linewidth]{./images/khd/gef-wind-rauch1.jpg}
\Captionof{figure}{\AuthNotesPicLegalOk{}{}Ein Auto brennt und alle schauen
faziniert zu \ldots \SourcePicture{Gabriel}} 
&
\includegraphics[width=\linewidth]{./images/khd/gef-wind-rauch2.jpg}
\Captionof{figure}{\AuthNotesPicLegalOk{}{}\ldots doch dann kommt Wind auf \ldots  
\SourcePicture{Gabriel}} 
&
\includegraphics[width=\linewidth]{./images/khd/gef-wind-rauch.jpg}
\Captionof{figure}{\AuthNotesPicLegalOk{}{}Und oh Schreck! Die Absperrung
verliert ihren Zweck!  \SourcePicture{Gabriel}} 
\\\bottomrule
\end{tabularx}

\Fazit{Die Moral aus der Geschicht': Vertrau' dem lauen Klima nicht!}
\end{gWideParEnv}
\end{table}

\TakeNotesLarge


\subsection{Spezialkräfte anfordern}
% TODO : Layout-Anpassung
\begin{anfxwarning}[author=GABS]{Layout-Anpassung notwendig.}
   Dieser Teil entspricht nicht dem Standard-Layout!
\end{anfxwarning}

\TakeNotesLarge
% TODO : Neufassung
\begin{anfxerror}[author=GABS]{Fehlender Text.}
  Text fehlt! Hier muß neuer Text geschrieben werden!
\end{anfxerror}

\nocite{Cicero911CommissionReport}